% !TeX root = ../main.tex
\chapter{PHƯƠNG PHÁP ĐỀ XUẤT}
\label{chap:phuongphap}

\section{Tổng quan quy trình}
\label{sec:tongquan}

Quy trình đề xuất gồm bốn giai đoạn nối tiếp, trong đó đầu ra của giai đoạn trước
trở thành đầu vào của giai đoạn sau. Điểm mấu chốt của thiết kế là \emph{không có
giai đoạn nào cần nhãn thủ công ở mức điểm ảnh}: toàn bộ thông tin không gian được
chiết xuất từ nhãn mức ảnh thông qua cơ chế giải thích của mạng phân loại.

\begin{figure}[H]
\centering
\resizebox{\textwidth}{!}{%
\begin{tikzpicture}[node distance=0.85cm and 0.95cm]
    \node[data] (dataset) {Bộ dữ liệu lá sầu riêng\\4.437 ảnh RGB\\5 lớp bệnh};
    \node[data, right=of dataset] (split) {Chia tập 70/10/20\\Train 3.104\\Val 443 / Test 890};
    \node[classmodel, right=of split] (classifier) {Giai đoạn 1\\EfficientNet-B0\\phân loại mức ảnh};
    \node[xai, right=of classifier] (gradcam) {Giai đoạn 2\\GradCAM++ đa tỉ lệ\\$0{,}6H_{\text{fine}}+0{,}4H_{\text{coarse}}$};
    \node[xai, right=of gradcam] (post) {Sinh mặt nạ thích ứng\\ngưỡng phân vị theo lớp\\hình thái + top-3 CC};
    \node[sam, right=of post] (samrefine) {Giai đoạn 3\\Tinh chỉnh bằng SAM\\gợi ý điểm FG/BG\\bảo vệ IoU (min = 0,15)};
    \node[segmodel, right=of samrefine] (segtrain) {Giai đoạn 4\\Phân đoạn UNet++\\bộ mã hoá EfficientNet-B0\\hàm mất mát FocalDice};

    \node[data, below=2.15cm of split] (newimage) {Ảnh lá đầu vào\\$224\times224\times3$};
    \node[classmodel, right=of newimage] (inferclass) {Dự đoán lớp bệnh\\kèm xác suất};
    \node[segmodel, right=of inferclass] (inferseg) {Định vị tổn thương\\mặt nạ nhị phân};
    \node[data, right=of inferseg] (decision) {Kết quả cuối\\lớp bệnh +\\vùng nhiễm bệnh};

    \draw[flow] (dataset) -- (split);
    \draw[flow] (split) -- (classifier);
    \draw[flow] (classifier) -- (gradcam);
    \draw[flow] (gradcam) -- (post);
    \draw[flow] (post) -- (samrefine);
    \draw[flow] (samrefine) -- (segtrain);

    \draw[dashedflow] (classifier.south) -- ++(0,-0.7) -| (inferclass.north);
    \draw[dashedflow] (segtrain.south) -- ++(0,-0.7) -| (inferseg.north);
    \draw[flow] (newimage) -- (inferclass);
    \draw[flow] (inferclass) -- (inferseg);
    \draw[flow] (inferseg) -- (decision);

    \node[lane, fit=(dataset)(split)(classifier)(gradcam)(post)(samrefine)(segtrain), inner xsep=7pt, inner ysep=12pt,
        label={[font=\bfseries\small]above:Huấn luyện và xây dựng nhãn giả}] {};
    \node[lane, fit=(newimage)(inferclass)(inferseg)(decision), inner xsep=7pt, inner ysep=12pt,
        label={[font=\bfseries\small]below:Suy luận}] {};
\end{tikzpicture}%
}
\caption[Tổng quan quy trình học giám sát yếu đề xuất]{Tổng quan quy trình phân
tích bệnh lá sầu riêng theo hướng giám sát yếu. Mô hình phân loại học nhãn mức ảnh
trước, sau đó các giải thích GradCAM++ đa tỉ lệ của nó được chuyển thành mặt nạ
phân đoạn giả thông qua ngưỡng phân vị theo lớp. SAM tinh chỉnh biên của các mặt
nạ này trước khi chúng được dùng để huấn luyện mô hình phân đoạn UNet++ cuối cùng.
Khi suy luận, hệ thống trả về đồng thời lớp bệnh và vùng tổn thương được định vị.}
\label{fig:pipeline}
\end{figure}

Bảng dưới tóm tắt vai trò và sản phẩm của từng giai đoạn.

\begin{table}[H]
\centering
\caption{Vai trò và sản phẩm của bốn giai đoạn}
\label{tab:giaidoan}
\small
\begin{tabular}{L{1.4cm} L{3.6cm} L{3.4cm} L{4.2cm}}
\toprule
\textbf{Giai đoạn} & \textbf{Nhiệm vụ} & \textbf{Đầu vào} & \textbf{Sản phẩm} \\
\midrule
1 & Phân loại bệnh & Ảnh + nhãn mức ảnh & Mô hình phân loại đã huấn luyện \\
2 & Sinh nhãn giả & Mô hình phân loại + 3.104 ảnh huấn luyện & 3.104 mặt nạ nhị phân \\
3 & Tinh chỉnh biên & Mặt nạ giả + ảnh gốc & 3.104 mặt nạ biên sắc nét hơn \\
4 & Phân đoạn & Ảnh + mặt nạ giả & Mô hình phân đoạn đầu cuối \\
\bottomrule
\end{tabular}
\end{table}

Luận văn xây dựng \emph{ba phiên bản} của quy trình, khác nhau ở cách sinh nhãn
giả ở Giai đoạn 2--3, còn Giai đoạn 4 giữ nguyên kiến trúc để bảo đảm tính so
sánh:

\begin{itemize}
    \item \textbf{V1} --- \GradCAM{} một tầng, ngưỡng cố định 0,5
          (\cref{sec:pseudov1}).
    \item \textbf{V2} --- \GradCAMpp{} đa tỉ lệ, ngưỡng phân vị theo lớp
          (\cref{sec:pseudov2}).
    \item \textbf{V3} --- nhãn V2 được \SAM{} tinh chỉnh biên
          (\cref{sec:samrefine}).
\end{itemize}

\section{Bộ dữ liệu và phân tích khám phá}
\label{sec:dulieu}

Phân tích khám phá dữ liệu (\en{Exploratory Data Analysis}, EDA) được thực hiện
trước khi huấn luyện nhằm ba mục đích: xác nhận tính toàn vẹn của dữ liệu, phát
hiện sớm các đặc điểm có thể gây khó khăn cho mô hình, và định hướng lựa chọn chỉ
số đánh giá.

\subsection{Thành phần và phân bố lớp}

Bộ dữ liệu gồm \num{4437} ảnh chia thành năm lớp như trong \cref{tab:phanbolop}.

\begin{table}[H]
\centering
\caption{Phân bố số lượng ảnh theo lớp}
\label{tab:phanbolop}
\begin{tabular}{c l r r}
\toprule
\textbf{Nhãn} & \textbf{Lớp} & \textbf{Số ảnh} & \textbf{Tỉ lệ} \\
\midrule
0 & Algal Leaf Spot (đốm tảo)          & 733 & 16,5\% \\
1 & Allocaridara Attack (rầy phấn)     & 913 & 20,6\% \\
2 & Healthy Leaf (lá khoẻ)             & 976 & 22,0\% \\
3 & Leaf Blight (cháy lá)              & 937 & 21,1\% \\
4 & Phomopsis Leaf Spot (đốm Phomopsis)& 878 & 19,8\% \\
\midrule
\multicolumn{2}{l}{\textbf{Tổng}} & \textbf{4.437} & \textbf{100\%} \\
\bottomrule
\end{tabular}
\end{table}

Nhãn số được gán theo thứ tự bảng chữ cái của tên thư mục, thống nhất giữa mọi
tập con và mọi lần chạy.

\begin{figure}[H]
\centering
\includegraphics[width=\textwidth]{figures/eda_class_distribution.png}
\caption[Phân bố lớp của bộ dữ liệu]{Phân bố số lượng ảnh theo lớp (trái) và tỉ lệ
phần trăm (phải). Đường nét đứt biểu thị giá trị trung bình 887 ảnh mỗi lớp.}
\label{fig:class_distribution}
\end{figure}

Tỉ lệ mất cân bằng --- thương của lớp nhiều nhất trên lớp ít nhất --- bằng
$976/733 = 1{,}33$. Đây là mức rất thấp; theo thông lệ, chỉ khi tỉ lệ vượt khoảng
1:10 mới cần các biện pháp can thiệp như lấy mẫu quá mức hay trọng số theo lớp.
Do đó luận văn \emph{không} áp dụng biện pháp nào và huấn luyện trực tiếp trên
phân bố gốc. Cân bằng gần hoàn hảo cũng khiến \en{macro-F1} trở thành chỉ số tổng
hợp phù hợp, vì không lớp nào đủ áp đảo để bóp méo giá trị trung bình.

\subsection{Phân chia tập dữ liệu}

Bộ dữ liệu được chia theo tỉ lệ 70/10/20 với phân tầng theo lớp.

\begin{table}[H]
\centering
\caption{Phân chia tập dữ liệu cho bài toán phân loại}
\label{tab:split}
\begin{tabular}{l r r L{6.0cm}}
\toprule
\textbf{Tập} & \textbf{Số ảnh} & \textbf{Tỉ lệ} & \textbf{Vai trò} \\
\midrule
Huấn luyện   & 3.104 & 70,0\% & Cập nhật trọng số \\
Kiểm định    & 443   & 10,0\% & Dừng sớm, chọn điểm kiểm tra, điều chỉnh tốc độ học \\
Kiểm thử     & 890   & 20,1\% & Đánh giá cuối, dùng đúng một lần \\
\midrule
\textbf{Tổng} & \textbf{4.437} & \textbf{100\%} & \\
\bottomrule
\end{tabular}
\end{table}

\begin{figure}[H]
\centering
\includegraphics[width=\textwidth]{figures/eda_dataset_splits.png}
\caption[Phân chia tập dữ liệu]{Số lượng ảnh trong ba tập con (trái) và phân bố
lớp trong từng tập (phải). Phân bố lớp đồng đều trên cả ba tập, không có hiện
tượng lệch khi chia dữ liệu.}
\label{fig:dataset_splits}
\end{figure}

Lựa chọn tỉ lệ này phản ánh một cân nhắc cụ thể. Tập kiểm định chỉ chiếm 10\%
(443 ảnh, trung bình khoảng 89 ảnh mỗi lớp) --- vừa đủ để ước lượng hiệu năng phục
vụ dừng sớm --- trong khi phần tiết kiệm được dồn cho tập kiểm thử 20\% (890 ảnh).
Tập kiểm thử lớn cho ước lượng hiệu năng cuối cùng có phương sai thấp hơn, điều
quan trọng khi cần báo cáo một con số duy nhất làm kết quả chính thức.

Phân bố theo lớp trong tập huấn luyện --- bộ ảnh sẽ được dùng để sinh nhãn giả ở
các giai đoạn sau --- như sau: Algal Leaf Spot 513 ảnh, Allocaridara Attack 639,
Healthy Leaf 683, Leaf Blight 655 và Phomopsis Leaf Spot 614.

\subsection{Đặc điểm hình thái của các lớp}

\begin{figure}[H]
\centering
\includegraphics[width=\textwidth]{figures/eda_samples_per_class.png}
\caption[Ảnh mẫu của năm lớp]{Hai ảnh mẫu cho mỗi lớp bệnh. Sự khác biệt hình thái
giữa Algal Leaf Spot và Phomopsis Leaf Spot là nhỏ nhất trong năm lớp.}
\label{fig:samples_per_class}
\end{figure}

\begin{table}[H]
\centering
\caption{Đặc điểm hình thái quan sát được của từng lớp}
\label{tab:hinhthai}
\small
\begin{tabular}{L{3.4cm} L{9.6cm}}
\toprule
\textbf{Lớp} & \textbf{Biểu hiện trên lá} \\
\midrule
Algal Leaf Spot & Đốm tròn hoặc bầu dục màu vàng--nâu--xám, rải rác trên bề mặt lá, viền tương đối rõ. \\
Allocaridara Attack & Vết cắn, đường kẻ và lỗ thủng không đều do côn trùng; kết cấu lá bị phá vỡ mạnh. \\
Healthy Leaf & Lá xanh đồng màu, bề mặt nhẵn, không có tổn thương hay đổi màu. \\
Leaf Blight & Mảng hoại tử lớn màu nâu--đen, thường lan từ rìa lá vào trong, ranh giới không rõ. \\
Phomopsis Leaf Spot & Đốm nhỏ hơn Algal, màu nâu tối viền vàng nhạt, thường tụ thành cụm. \\
\bottomrule
\end{tabular}
\end{table}

Quan sát quan trọng nhất từ \cref{fig:samples_per_class} là cặp \emph{Algal Leaf
Spot} và \emph{Phomopsis Leaf Spot} có hình thái gần nhau: cả hai đều tạo đốm nhỏ
trên nền lá xanh, khác biệt chủ yếu ở kích thước và mật độ đốm. Đây là cặp lớp dự
kiến gây nhầm lẫn nhiều nhất, và dự đoán này sẽ được kiểm chứng bằng ma trận nhầm
lẫn ở \cref{chap:thucnghiem}.

\subsection{Phân tích màu sắc}

\begin{figure}[H]
\centering
\includegraphics[width=\textwidth]{figures/eda_rgb_histogram.png}
\caption[Biểu đồ phân bố kênh màu RGB]{Phân bố ba kênh màu R, G, B cho từng lớp.
Healthy Leaf là lớp duy nhất có kênh Green dịch rõ về phía giá trị cao.}
\label{fig:rgb_histogram}
\end{figure}

\begin{figure}[H]
\centering
\includegraphics[width=\textwidth]{figures/eda_mean_rgb.png}
\caption[Giá trị trung bình các kênh màu theo lớp]{Giá trị trung bình của ba kênh
màu theo từng lớp, tính trên 30 ảnh mỗi lớp.}
\label{fig:mean_rgb}
\end{figure}

\begin{table}[H]
\centering
\caption{Giá trị trung bình các kênh màu theo lớp (mẫu 30 ảnh mỗi lớp)}
\label{tab:meanrgb}
\begin{tabular}{l r r r r}
\toprule
\textbf{Lớp} & \textbf{R} & \textbf{G} & \textbf{B} & \textbf{G $-$ R} \\
\midrule
Algal Leaf Spot     & 95,8  & 125,8 & 106,6 & 30,0 \\
Allocaridara Attack & 94,4  & 134,5 & 98,9  & 40,1 \\
Healthy Leaf        & 96,4  & 133,0 & 92,1  & 36,6 \\
Leaf Blight         & 103,5 & 137,8 & 107,2 & 34,3 \\
Phomopsis Leaf Spot & 97,2  & 138,6 & 88,5  & 41,4 \\
\bottomrule
\end{tabular}
\end{table}

Phân tích màu cho hai nhận định có ý nghĩa thực tiễn.

Thứ nhất, \emph{Algal Leaf Spot có hiệu $G - R$ thấp nhất} (30,0) --- dấu hiệu của
việc mô lá chuyển sang nâu đỏ, làm giảm ưu thế của kênh lục. Đây là đặc trưng màu
sắc phân biệt được lớp này khỏi các lớp còn lại.

Thứ hai, và quan trọng hơn về mặt phương pháp, \emph{các giá trị trung bình toàn
ảnh không đủ để phân tách năm lớp}: khoảng biến thiên của kênh lục giữa các lớp
chỉ từ 125,8 đến 138,6, quá hẹp so với độ lệch trong nội bộ từng lớp. Điều này bác
bỏ khả năng xây dựng một bộ phân loại đơn giản dựa trên đặc trưng màu thủ công, và
củng cố lựa chọn dùng CNN để học đặc trưng \emph{cục bộ} --- vốn là nơi thông tin
bệnh học thực sự nằm.

\subsection{Độ sáng và dung lượng tệp}

\begin{figure}[H]
\centering
\includegraphics[width=\textwidth]{figures/eda_filesize_brightness.png}
\caption[Dung lượng tệp và độ sáng]{Phân bố dung lượng tệp (trái) và phân bố độ
sáng trung bình theo lớp (phải).}
\label{fig:filesize_brightness}
\end{figure}

\begin{figure}[H]
\centering
\includegraphics[width=\textwidth]{figures/eda_brightness_kde.png}
\caption[Mật độ phân bố độ sáng theo lớp]{Ước lượng mật độ hạt nhân của độ sáng
trung bình cho từng lớp.}
\label{fig:brightness_kde}
\end{figure}

Dung lượng tệp tập trung quanh giá trị trung bình 11,36 KB với độ lệch chuẩn chỉ
0,80 KB (nhỏ nhất 8,91 KB, lớn nhất 15,07 KB). Độ đồng đều cao này xác nhận toàn
bộ ảnh đã đi qua cùng một quy trình nén và chuẩn hoá --- nhất quán với nhận định
về trạng thái tiền xử lý của dữ liệu ở \cref{subsec:nguongoc}.

Về hình học, kiểm tra trên mẫu 100 ảnh cho thấy chiều rộng, chiều cao và tỉ lệ
khung hình đều có độ lệch chuẩn bằng 0 --- mọi ảnh đúng $224 \times 224$ và vuông.
Hệ quả trực tiếp: bước thay đổi kích thước trong quy trình tiền xử lý trở thành
phép đồng nhất, và không có biến dạng hình học nào được đưa thêm vào dữ liệu.

\subsection{Kiểm chứng luồng nạp dữ liệu}

Trước khi huấn luyện, luồng nạp dữ liệu được kiểm chứng bằng cách lấy một lô mẫu
và đối chiếu các đại lượng thống kê. Kết quả: kích thước lô $[32, 3, 224, 224]$
đúng như thiết kế; giá trị điểm ảnh sau chuẩn hoá nằm trong khoảng
$[-2{,}118;\ 2{,}640]$; nhãn nằm trong $[0, 4]$; độ lệch chuẩn theo từng kênh xấp
xỉ 1,0 (\num{1.0298}, \num{1.0201}, \num{1.0387}).

Giá trị trung bình theo kênh sau chuẩn hoá là $(-0{,}618;\ 0{,}001;\ -0{,}320)$
thay vì $(0, 0, 0)$. Đây \emph{không} phải lỗi mà là hệ quả có thể dự đoán được:
phép chuẩn hoá dùng giá trị trung bình và độ lệch chuẩn của ImageNet, trong khi
ảnh lá cây có kênh lục trội hơn và kênh đỏ, lam yếu hơn so với phân bố ImageNet.
Độ lệch chuẩn theo kênh xấp xỉ 1,0 cho thấy phép chuẩn hoá vẫn phát huy đúng tác
dụng chính là đưa dữ liệu về thang đo phù hợp cho tối ưu bằng gradient.

\section{Giai đoạn 1 --- Mô hình phân loại}
\label{sec:giaidoan1}

\subsection{Kiến trúc}

Mô hình phân loại dùng EfficientNet-B0 tiền huấn luyện trên ImageNet, thay tầng
phân loại 1000 lớp gốc bằng cấu trúc:

\begin{equation}
    \text{Dropout}(p = 0{,}2) \;\rightarrow\; \text{Linear}(1280 \rightarrow 5)
    \label{eq:cls_head}
\end{equation}

Toàn bộ mô hình có \num{4013953} tham số và được tinh chỉnh hoàn toàn --- không
tầng nào bị đóng băng --- theo biện luận ở \cref{sec:cnn}.

\subsection{Tăng cường dữ liệu}

Tăng cường dữ liệu (\en{data augmentation}) được áp dụng \emph{chỉ trên tập huấn
luyện}: lật ngang, lật dọc, xoay và biến đổi màu (\en{color jitter}). Tập kiểm
định và kiểm thử không được tăng cường, vì nếu tăng cường thì mỗi lần đánh giá sẽ
cho kết quả khác nhau, làm nhiễu chính tín hiệu dùng để dừng sớm.

Lật dọc là phép biến đổi hợp lệ trong bài toán này --- khác với ảnh chữ viết hay
ảnh phong cảnh --- vì một chiếc lá không có hướng ``đúng'' cố định khi chụp.

\subsection{Cấu hình huấn luyện}

\begin{table}[H]
\centering
\caption{Cấu hình huấn luyện mô hình phân loại}
\label{tab:cauhinh_cls}
\begin{tabular}{L{4.6cm} L{8.4cm}}
\toprule
\textbf{Thành phần} & \textbf{Thiết lập} \\
\midrule
Kiến trúc          & EfficientNet-B0, tiền huấn luyện ImageNet \\
Tầng phân loại     & Dropout(0,2) + Linear(1280 $\rightarrow$ 5) \\
Số tham số         & 4.013.953 (toàn bộ khả huấn luyện) \\
Hàm mất mát        & Cross-Entropy \\
Bộ tối ưu          & Adam, tốc độ học $1 \times 10^{-4}$ \\
Lịch tốc độ học    & ReduceLROnPlateau theo \texttt{val\_acc}, patience = 2, hệ số 0,5 \\
Kích thước lô      & 16 \\
Số epoch tối đa    & 50 \\
Dừng sớm           & patience = 5 theo \texttt{val\_acc} \\
Tăng cường         & Lật ngang/dọc, xoay, biến đổi màu \\
\bottomrule
\end{tabular}
\end{table}

Hai lựa chọn cần biện luận:

\textbf{Tốc độ học $10^{-4}$.} Giá trị này nhỏ hơn một bậc so với mặc định
$10^{-3}$ của Adam. Với tinh chỉnh toàn bộ, tốc độ học lớn khiến các bước cập nhật
đầu tiên --- khi tầng phân loại mới còn khởi tạo ngẫu nhiên và sinh gradient lớn
--- phá huỷ chính các trọng số tiền huấn luyện mà ta muốn tận dụng.

\textbf{Kích thước lô 16.} Lô nhỏ tạo nhiễu gradient, mà nhiễu này có tác dụng
chính quy hoá nhẹ --- hữu ích với tập huấn luyện chỉ hơn ba nghìn ảnh.

\subsection{Chiến lược dừng sớm}

Cần phân biệt rõ hai cơ chế cùng theo dõi \texttt{val\_acc} nhưng phục vụ hai mục
đích khác nhau: \en{ReduceLROnPlateau} với patience = 2 \emph{giảm tốc độ học} khi
hiệu năng chững lại, còn dừng sớm với patience = 5 \emph{kết thúc huấn luyện}. Việc
đặt patience của bộ lập lịch nhỏ hơn của dừng sớm là có chủ đích: mô hình được cho
ít nhất hai cơ hội giảm tốc độ học để thoát vùng phẳng trước khi bị dừng hẳn.

Điểm kiểm tra chỉ được ghi đè khi \texttt{val\_acc} vượt kỷ lục cũ, nên tệp
\texttt{best\_\allowbreak model.\allowbreak pth} luôn ứng với mô hình tốt nhất từng đạt được. Đây chính là
tệp được nạp lại ở tất cả các giai đoạn sau.

\section{Giai đoạn 2 --- Sinh nhãn giả phiên bản V1}
\label{sec:pseudov1}

\subsection{Trích xuất bản đồ nhiệt}

Nhãn giả V1 dùng \GradCAM{} một tầng theo \cref{eq:gradcam}, với tầng đích là
\texttt{base.\allowbreak features.\allowbreak 8.\allowbreak 0} --- tầng tích chập cuối cùng của EfficientNet-B0, cho
bản đồ đặc trưng $7 \times 7$. Đây là lựa chọn tiêu chuẩn: tầng sâu nhất giữ nhiều
ngữ nghĩa nhất trong khi vẫn còn thông tin không gian.

Quy trình được áp dụng cho \emph{toàn bộ} \num{3104} ảnh huấn luyện. Phép biến đổi
đầu vào giống hệt tập kiểm định và kiểm thử, không tăng cường --- yêu cầu bắt buộc
để mặt nạ sinh ra ổn định giữa các lần chạy.

\subsection{Nhị phân hoá và hậu xử lý hình thái}

Bản đồ nhiệt đã chuẩn hoá về $[0, 1]$ được nhị phân hoá bằng ngưỡng cố định:

\begin{equation}
    M(i,j) = \begin{cases}
        1 & \text{nếu } H(i,j) \geq 0{,}5 \\
        0 & \text{nếu ngược lại}
    \end{cases}
    \label{eq:threshold_v1}
\end{equation}

Sau đó áp dụng lần lượt phép \emph{đóng} rồi phép \emph{mở} theo
\cref{eq:closing,eq:opening} với phần tử cấu trúc \emph{chữ nhật} $5 \times 5$.
Kích thước này tương ứng khoảng 0,05\% diện tích ảnh --- đủ để làm mịn biên mà
không xoá mất các đốm bệnh nhỏ.

\begin{figure}[H]
\centering
\includegraphics[width=\textwidth]{figures/pl1_morphology_demo.png}
\caption[Minh hoạ hậu xử lý hình thái]{Tác động của hậu xử lý hình thái. Từ trái
sang phải: ảnh gốc, bản đồ nhiệt Grad-CAM, mặt nạ chỉ qua ngưỡng hoá, và mặt nạ
sau khi áp dụng phép đóng rồi phép mở với phần tử cấu trúc $5 \times 5$.}
\label{fig:morphology_demo}
\end{figure}

Toàn bộ quá trình sinh \num{3104} mặt nạ mất 14 phút 00 giây trên GPU, tương đương
tốc độ 3,69 ảnh mỗi giây.

\subsection{Ba hạn chế của phiên bản V1}
\label{subsec:hanche_v1}

Phân tích thống kê mặt nạ V1 bộc lộ ba vấn đề, mỗi vấn đề dẫn tới một cải tiến cụ
thể ở phiên bản V2.

\textbf{Thứ nhất --- lớp lá khoẻ sinh mặt nạ vô nghĩa.} Độ phủ trung bình của lớp
Healthy Leaf đạt 26,4\%, cao gần gấp đôi các lớp bệnh (12,3--15,1\%), với giá trị
lớn nhất lên tới 59,1\%. Điều này hoàn toàn logic mà lại rất tai hại: lá khoẻ
\emph{không có} tổn thương nào để mạng tập trung, nên bản đồ chú ý khuếch tán khắp
phiến lá; ngưỡng hoá một bản đồ khuếch tán tất yếu cho vùng rộng. Nếu đưa
\num{683} mặt nạ này vào huấn luyện phân đoạn, mô hình sẽ học rằng lá khoẻ cũng có
26\% diện tích ``bệnh''.

\textbf{Thứ hai --- ngưỡng cố định không thích ứng.} Ngưỡng 0,5 áp dụng cho mọi
ảnh, bất kể phân bố bản đồ nhiệt của từng ảnh. Ảnh bệnh nhẹ bị phân đoạn thiếu,
ảnh có nền nổi bật bị phân đoạn thừa. Hệ quả là độ lệch chuẩn của độ phủ rất lớn
(0,039--0,076 tuỳ lớp).

\textbf{Thứ ba --- biên quá thô.} Bản đồ $7 \times 7$ phóng lên $224 \times 224$
khiến mỗi ô tương ứng một khối $32 \times 32$ điểm ảnh. Các đốm bệnh rời rạc bị
gộp thành một khối lớn duy nhất, làm mất thông tin về số lượng và vị trí từng đốm
--- đúng như phân tích lý thuyết ở \cref{sec:xai}.

\section{Giai đoạn 2' --- Sinh nhãn giả phiên bản V2}
\label{sec:pseudov2}

Phiên bản V2 khắc phục cả ba hạn chế trên bằng bốn thay đổi kỹ thuật, tóm tắt ở
\cref{tab:v1_v2}.

\begin{table}[H]
\centering
\caption{So sánh thiết kế nhãn giả V1 và V2}
\label{tab:v1_v2}
\small
\begin{tabular}{L{2.6cm} L{3.5cm} L{3.5cm} L{3.6cm}}
\toprule
\textbf{Thành phần} & \textbf{V1} & \textbf{V2} & \textbf{Lý do thay đổi} \\
\midrule
Phương pháp CAM & Grad-CAM & Grad-CAM++ & Tránh hiệu ứng san phẳng khi có nhiều đốm bệnh \\
Tầng đích & \texttt{features.8} ($7\times7$) & \texttt{features.6} ($14\times14$) + \texttt{features.8} & Tăng chi tiết không gian gấp bốn lần về diện tích \\
Ngưỡng hoá & Cố định 0,5 & Phân vị theo lớp & Kiểm soát trực tiếp độ phủ \\
Phần tử cấu trúc & Chữ nhật $5\times5$ & Elip $5\times5$ & Bám hình dạng tròn của đốm bệnh \\
Lọc nhiễu & Không & Top-3 CC $\geq$ 400 px & Loại mảnh vụn mà phép mở bỏ sót \\
Lá khoẻ & Xử lý như lớp bệnh & Mặt nạ toàn số 0 & Loại 683 mặt nạ dương tính giả \\
\bottomrule
\end{tabular}
\end{table}

\subsection{Grad-CAM++ đa tỉ lệ}

Hai tầng đích được chọn theo \cref{eq:multiscale}:

\begin{itemize}
    \item \textbf{Tầng tinh}: \texttt{base.\allowbreak features.\allowbreak 6.\allowbreak 3.\allowbreak block.\allowbreak 3.\allowbreak 0}, độ phân giải
          $14 \times 14$ --- nắm bắt các đốm nhỏ.
    \item \textbf{Tầng thô}: \texttt{base.\allowbreak features.\allowbreak 8.\allowbreak 0}, độ phân giải
          $7 \times 7$ --- cung cấp ngữ cảnh ngữ nghĩa.
\end{itemize}

Bản đồ nhiệt hợp nhất:

\begin{equation}
    H_{\text{final}} = 0{,}6 \cdot H_{\text{fine}} + 0{,}4 \cdot H_{\text{coarse}}
    \label{eq:fusion_v2}
\end{equation}

Trọng số thiên về tầng tinh nhằm ưu tiên độ sắc nét của biên, nhưng vẫn giữ 40\%
đóng góp từ tầng thô để tránh kích hoạt nhầm trên kết cấu nền. Trọng số của từng
tầng được tính theo \cref{eq:gradcampp_alpha,eq:gradcampp_w}.

\subsection{Ngưỡng phân vị theo lớp}

Đây là cải tiến trung tâm của V2. Thay vì cố định giá trị ngưỡng, ta cố định
\emph{tỉ lệ diện tích} được giữ lại. Với tỉ lệ mục tiêu $p$, ngưỡng được tính là
phân vị thứ $(1 - p) \times 100$ của chính bản đồ nhiệt đó:

\begin{equation}
    \tau = \operatorname{percentile}\big(H,\; (1 - p) \times 100\big),
    \qquad
    M(i,j) = \mathbf{1}\big[H(i,j) \geq \tau\big]
    \label{eq:percentile}
\end{equation}

Cách làm này đảo ngược quan hệ nhân quả so với V1: ở V1 ngưỡng là tham số và độ
phủ là kết quả không kiểm soát được; ở V2 độ phủ là tham số còn ngưỡng được suy ra
cho từng ảnh. Nhờ đó ngưỡng tự thích ứng với phân bố riêng của từng bản đồ nhiệt.

\begin{table}[H]
\centering
\caption{Tỉ lệ độ phủ mục tiêu theo lớp}
\label{tab:keeppct}
\begin{tabular}{l c L{7.2cm}}
\toprule
\textbf{Lớp} & \textbf{Giữ lại} & \textbf{Căn cứ} \\
\midrule
Algal Leaf Spot     & 18\% & Đốm nhỏ rải rác, cần độ chính xác cao \\
Allocaridara Attack & 20\% & Tổn thương côn trùng trên vùng rộng hơn \\
Healthy Leaf        & 0\%  & Không có mô tổn thương --- ép về không \\
Leaf Blight         & 22\% & Mảng hoại tử lan rộng nhất \\
Phomopsis Leaf Spot & 18\% & Đốm nhỏ, tương tự Algal \\
\bottomrule
\end{tabular}
\end{table}

Các giá trị 18--22\% được chọn dựa trên đặc điểm bệnh học quan sát từ
\cref{tab:hinhthai}, theo thứ tự: bệnh đốm chiếm diện tích nhỏ nhất, bệnh cháy lá
lan rộng nhất. Cần nói rõ đây là \emph{tham số suy luận, chưa được hiệu chỉnh dựa
trên nhãn chuẩn} --- một hạn chế đã nêu ở \cref{subsec:gioihan-quantrong}.

Lý do không dùng ngưỡng Otsu đã được phân tích ở \cref{sec:xulyanh}: phép hợp nhất
đa tỉ lệ theo \cref{eq:fusion_v2} làm bản đồ nhiệt trở nên trơn và đơn đỉnh, vi
phạm giả định hai đỉnh mà Otsu dựa vào.

\subsection{Ép về không cho lớp lá khoẻ}

Với lớp Healthy Leaf, tỉ lệ giữ lại được đặt bằng 0, tức mặt nạ là ma trận không:

\begin{equation}
    M_{\text{healthy}}(i,j) = 0 \quad \forall (i,j)
    \label{eq:hardzero}
\end{equation}

Đây là một \emph{tri thức tiên nghiệm} được đưa thẳng vào quy trình: lá khoẻ theo
định nghĩa không có tổn thương, nên mọi kích hoạt trên ảnh lá khoẻ đều là dương
tính giả xét theo mục tiêu phân đoạn bệnh. Biện pháp này loại bỏ \num{683} mặt nạ
nhiễu và cung cấp cho mô hình phân đoạn các mẫu âm sạch --- điều mà bản thân bản đồ
kích hoạt không bao giờ tự cho được.

\subsection{Hậu xử lý}

Hai bước nối tiếp sau ngưỡng hoá:

\begin{enumerate}
    \item \textbf{Hình thái với phần tử elip $5 \times 5$}: đóng rồi mở. Phần tử
          elip bám sát hình dạng tròn hoặc bầu dục của đốm bệnh tốt hơn phần tử
          chữ nhật dùng ở V1.
    \item \textbf{Lọc top-3 thành phần liên thông}, chỉ giữ các vùng có diện tích
          $\geq 400$ điểm ảnh.
\end{enumerate}

Hai tham số của bước lọc phản ánh đặc thù bài toán. Giữ \emph{ba} vùng thay vì một
vì bệnh đốm lá biểu hiện thành nhiều tổn thương tách biệt --- giữ một vùng sẽ làm
mất phần lớn thông tin. Ngưỡng \emph{400 điểm ảnh}, tương đương khoảng 0,8\% diện
tích ảnh, đủ lớn để loại các mảnh vụn cỡ trung mà phép mở hình thái không xử lý
hết, nhưng vẫn nhỏ hơn một đốm bệnh điển hình.

\begin{figure}[H]
\centering
\includegraphics[width=\textwidth]{figures/pl2_v1_vs_v2.png}
\caption[So sánh mặt nạ V1 và V2]{So sánh chất lượng mặt nạ trên cùng một ảnh. Từ
trái sang phải: ảnh gốc, bản đồ nhiệt GradCAM++ đa tỉ lệ, mặt nạ V1 (ngưỡng cố
định 0,5, phần tử chữ nhật) và mặt nạ V2 (ngưỡng phân vị, phần tử elip, lọc top-3
thành phần liên thông).}
\label{fig:v1_vs_v2}
\end{figure}

Toàn bộ quá trình sinh nhãn V2 mất 4 phút 08 giây (12,51 ảnh mỗi giây) --- nhanh
hơn V1 khoảng 3,4 lần dù tính hai bản đồ kích hoạt thay vì một, nhờ cách cài đặt
hiệu quả hơn cho phép truyền ngược.

\subsection{Đánh đổi của thiết kế V2}
\label{subsec:tradeoff_v2}

Ngưỡng phân vị bảo đảm độ phủ bám sát mục tiêu, nhưng cái giá phải trả cần được
nêu thẳng: \emph{độ phủ trở nên gần như đồng nhất trong cùng một lớp}. Một ảnh chỉ
có một đốm bệnh nhỏ và một ảnh bị bệnh phủ kín lá đều nhận mặt nạ chiếm xấp xỉ
cùng tỉ lệ diện tích.

Nói cách khác, mặt nạ V2 \emph{không mã hoá được mức độ nặng nhẹ của bệnh}. Mô
hình phân đoạn huấn luyện trên nhãn này học cách ``phủ đúng diện tích'' hơn là học
cách ``phát hiện vùng bệnh thực sự''. Đây là đánh đổi được chấp nhận có ý thức để
loại bỏ hiện tượng phân đoạn thừa của Otsu, và hệ quả của nó sẽ hiện rõ trong kết
quả thực nghiệm ở \cref{chap:thucnghiem}.

\section{Giai đoạn 3 --- Tinh chỉnh biên bằng SAM}
\label{sec:samrefine}

\subsection{Động lực}

Mặt nạ V2 định vị đúng vùng bệnh nhưng biên vẫn thô, vì bản đồ nhiệt gốc chỉ có độ
phân giải $14 \times 14$ và $7 \times 7$. \SAM{} được kỳ vọng bổ khuyết đúng điểm
này: giữ nguyên thông tin \emph{vị trí} từ V2 và thay thế thông tin \emph{biên}
bằng biên chính xác ở mức điểm ảnh mà \SAM{} sinh ra.

\subsection{Sinh gợi ý}

Với mỗi mặt nạ V2 khác rỗng, các gợi ý được trích xuất như sau:

\begin{itemize}
    \item \textbf{Điểm tiền cảnh} (5 điểm): trọng tâm của mặt nạ, cộng 4 điểm lấy
          ngẫu nhiên trong vùng mặt nạ. Việc bổ sung điểm ngẫu nhiên là cần thiết
          --- nếu chỉ dùng trọng tâm thì với mặt nạ hình vành khuyên hoặc gồm
          nhiều cụm rời, trọng tâm có thể rơi ra ngoài vùng thực.
    \item \textbf{Điểm nền} (3 điểm): lấy ngẫu nhiên trong phần bù của mặt nạ đã
          giãn nở 30 điểm ảnh. Vùng đệm 30 điểm ảnh bảo đảm điểm nền không rơi sát
          biên tổn thương --- nơi nhãn vốn không chắc chắn.
\end{itemize}

Mặt nạ rỗng --- toàn bộ \num{683} ảnh lá khoẻ --- được bỏ qua và giữ nguyên trạng
thái rỗng.

\subsection{Cơ chế bảo vệ dựa trên IoU}

Để phòng trường hợp \SAM{} chọn nhầm vùng, một cơ chế bảo vệ được cài đặt:

\begin{equation}
    M_{\text{refined}} =
    \begin{cases}
        M_{\SAM} & \text{nếu } \operatorname{IoU}(M_{\SAM},\, M_{\text{V2}}) \geq 0{,}15 \\[4pt]
        M_{\text{V2}} & \text{nếu ngược lại}
    \end{cases}
    \label{eq:iou_guard}
\end{equation}

Ý tưởng là nếu mặt nạ \SAM{} sinh ra chồng lấp quá ít với mặt nạ V2 thì nhiều khả
năng \SAM{} đã chuyển sang một đối tượng khác, khi đó ta quay về dùng mặt nạ V2.

Cần lưu ý ngay rằng cơ chế này \emph{chỉ chặn được một trong hai kiểu sai lệch}.
Phân tích đầy đủ về khiếm khuyết của nó, cùng bằng chứng định lượng, được trình
bày ở \cref{chap:thucnghiem} --- đây là một trong những phát hiện chính của luận
văn.

\subsection{Hậu xử lý}

Mặt nạ sau tinh chỉnh được làm sạch bằng hai bước: loại các thành phần liên thông
có diện tích nhỏ hơn 5\% tổng diện tích mặt nạ hoặc nhỏ hơn 30 điểm ảnh, sau đó
áp dụng phép đóng hình thái $5 \times 5$ để lấp các lỗ nhỏ.

\section{Giai đoạn 4 --- Mô hình phân đoạn}
\label{sec:giaidoan4}

\subsection{Kiến trúc}

Phiên bản V1 dùng EfficientNet-UNet tự cài đặt: bộ mã hoá EfficientNet-B0 trích
đặc trưng tại bốn mức, bộ giải mã U-Net với kết nối tắt trực tiếp, tổng cộng
\num{7276925} tham số (bộ mã hoá \num{4007548}, bộ giải mã \num{3269312}).

Hai phiên bản V2 và V3 dùng \UNetpp{} với bộ mã hoá EfficientNet-B0
(\num{6569581} tham số) theo biện luận ở \cref{sec:unet}. Đầu ra là một kênh
\en{logit}, hàm \en{sigmoid} được đưa vào trong hàm mất mát thay vì đặt ở đầu ra
mô hình --- cách làm chuẩn để bảo đảm ổn định số học.

\subsection{Tăng cường dữ liệu cho phân đoạn}

Khác với phân loại, tăng cường cho phân đoạn phải biến đổi \emph{đồng bộ} cả ảnh
lẫn mặt nạ. Phiên bản V1 chỉ dùng lật ngang và lật dọc. Hai phiên bản V2 và V3 dùng
bộ tăng cường mạnh hơn nhiều qua thư viện
Albumentations~\cite{buslaev2020albumentations}, liệt kê ở \cref{tab:augment}.

\begin{table}[H]
\centering
\caption{Bộ tăng cường dữ liệu cho phân đoạn (V2 và V3)}
\label{tab:augment}
\small
\begin{tabular}{L{4.2cm} c L{7.2cm}}
\toprule
\textbf{Phép biến đổi} & \textbf{Xác suất} & \textbf{Mục đích} \\
\midrule
RandomResizedCrop (0,7--1,0) & 1,00 & Bất biến với tỉ lệ và vị trí tổn thương \\
HorizontalFlip / VerticalFlip & 0,50 & Bất biến với hướng đặt lá \\
RandomRotate90 & 0,50 & Bất biến với góc quay vuông \\
ShiftScaleRotate & 0,50 & Bất biến với xoay nhỏ và tịnh tiến \\
ElasticTransform & 0,30 & Mô phỏng lá cong vênh trong thực tế \\
GridDistortion & 0,20 & Mô phỏng biến dạng phối cảnh \\
ColorJitter & 0,50 & Bất biến với điều kiện ánh sáng \\
GaussNoise & 0,20 & Bền vững với nhiễu cảm biến \\
\bottomrule
\end{tabular}
\end{table}

Hai phép biến đổi đàn hồi và biến dạng lưới đáng chú ý vì chúng làm biến dạng
\emph{cả mặt nạ}. Với nhãn chuẩn, điều này chấp nhận được; với nhãn giả vốn đã
không chắc chắn về biên, việc biến dạng thêm có thể làm nhiễu tín hiệu huấn luyện.
Đây là một điểm cần cân nhắc, được thảo luận lại ở \cref{chap:thucnghiem}.

\subsection{Cấu hình của ba phiên bản}

\begin{table}[H]
\centering
\caption{Cấu hình huấn luyện ba phiên bản mô hình phân đoạn}
\label{tab:cauhinh_seg}
\small
\begin{tabular}{L{3.2cm} L{3.2cm} L{3.2cm} L{3.2cm}}
\toprule
\textbf{Thành phần} & \textbf{V1} & \textbf{V2} & \textbf{V3} \\
\midrule
Kiến trúc & EfficientNet-UNet & UNet++ EB0 & UNet++ EB0 \\
Số tham số & 7.276.925 & 6.569.581 & 6.569.581 \\
Nhãn huấn luyện & Pseudo V1 & Pseudo V2 & SAM tinh chỉnh \\
Hàm mất mát & BCE + Dice & FocalDice & FocalDice \\
Bộ tối ưu & Adam $10^{-4}$ & AdamW $10^{-4}$, wd $10^{-4}$ & AdamW $10^{-4}$, wd $10^{-4}$ \\
Lịch tốc độ học & ReduceLROnPlateau & CosineAnnealing $T_{\max}=50$ & CosineAnnealing $T_{\max}=50$ \\
Kích thước lô & 8 & 4 & 8 \\
Epoch tối đa & 30 & 50 & 50 \\
Dừng sớm & patience 7 & patience 5 & patience 10 \\
Tiêu chí dừng & \texttt{val\_loss} & \texttt{val\_loss} & \texttt{val\_loss} \\
Cắt gradient & Không & \texttt{max\_norm} = 1,0 & \texttt{max\_norm} = 1,0 \\
Tăng cường & Chỉ lật & Albumentations & Albumentations \\
\bottomrule
\end{tabular}
\end{table}

Thiết kế thực nghiệm được tổ chức sao cho \textbf{V2 và V3 chỉ khác nhau ở một
biến duy nhất là bộ nhãn huấn luyện}. Kiến trúc, hàm mất mát, bộ tối ưu, lịch tốc
độ học, bộ tăng cường và cách chia tập đều giữ nguyên. Nhờ vậy, chênh lệch hiệu
năng giữa hai phiên bản có thể quy cho tác động của việc tinh chỉnh biên bằng
\SAM{} --- đây là phép so sánh có kiểm soát duy nhất trong luận văn.

Hai khác biệt còn lại giữa V2 và V3 --- kích thước lô 4 so với 8 và patience 5 so
với 10 --- xuất phát từ ràng buộc kỹ thuật khi chạy thực nghiệm chứ không phải lựa
chọn thiết kế, và ảnh hưởng của chúng được đánh giá ở \cref{chap:thucnghiem}.

\subsection{Chia tập cho bài toán phân đoạn}
\label{subsec:split_seg}

Toàn bộ \num{3104} ảnh huấn luyện kèm nhãn giả được chia lại theo tỉ lệ 70/15/15,
cho \num{2172} ảnh huấn luyện, 465 ảnh kiểm định và 467 ảnh kiểm thử, với hạt
giống ngẫu nhiên cố định bằng 42.

Cần lưu ý một chi tiết kỹ thuật có hệ quả trực tiếp tới việc diễn giải kết quả:

\begin{quote}
Phiên bản V1 chia tập bằng \texttt{torch.\allowbreak random\_\allowbreak split} với bộ sinh số ngẫu nhiên
của PyTorch, trong khi V2 và V3 chia bằng
\texttt{numpy.\allowbreak random.\allowbreak default\_rng(42).\allowbreak
permutation}. Dù cùng hạt giống 42, hai bộ
sinh số ngẫu nhiên này tạo ra \emph{hai hoán vị khác nhau}, nên \textbf{tập kiểm
thử của V1 không trùng với tập kiểm thử của V2 và V3}.
\end{quote}

Đây là một trong nhiều lý do khiến chỉ số của V1 không so sánh trực tiếp được với
V2 và V3, bên cạnh lý do quan trọng hơn về khác biệt mốc tham chiếu, được phân
tích trọn vẹn ở \cref{sec:caveat}. Ngược lại, V2 và V3 dùng chung đúng một hoán vị
nên tập kiểm thử của chúng trùng khớp hoàn toàn.

\section{Kết luận chương}

Chương này đã trình bày chi tiết quy trình bốn giai đoạn cùng ba phiên bản nhãn
giả. Các quyết định thiết kế chính và căn cứ của chúng được tóm tắt như sau:

\begin{itemize}
    \item Dùng chung EfficientNet-B0 cho cả phân loại và bộ mã hoá phân đoạn, bảo
          đảm bản đồ kích hoạt và đặc trưng phân đoạn cùng không gian biểu diễn.
    \item Chuyển từ \GradCAM{} một tầng sang \GradCAMpp{} đa tỉ lệ để xử lý đặc
          thù nhiều đốm bệnh trên một lá.
    \item Thay ngưỡng cố định bằng ngưỡng phân vị theo lớp, đảo ngược quan hệ giữa
          tham số và kết quả để kiểm soát trực tiếp độ phủ.
    \item Ép mặt nạ lớp lá khoẻ về không, đưa tri thức tiên nghiệm vào quy trình
          để loại \num{683} mặt nạ dương tính giả.
    \item Giữ nguyên mọi yếu tố giữa V2 và V3 ngoại trừ bộ nhãn, tạo ra phép so
          sánh có kiểm soát cho tác động của \SAM{}.
\end{itemize}

Hai điểm đã được nêu trước và sẽ được kiểm chứng bằng số liệu ở chương sau: đánh
đổi giữa kiểm soát độ phủ và khả năng mã hoá mức độ bệnh
(\cref{subsec:tradeoff_v2}), và khiếm khuyết của cơ chế bảo vệ dựa trên IoU
(\cref{eq:iou_guard}).
